% !PRINTED 450
\weberpage{450}
\begin{center}
§ 123. \gesperrt{Beziehungen zwischen Klasseninvarianten in den}\\
\gesperrt{verschiedenen Ordnungen.}
\end{center}

Zwischen den Klasseninvarianten verschiedener Diskriminanten $D = Q^2\Delta$ mit dem gleichen Stamm $\Delta$ bestehen algebraische Beziehungen, die wir jetzt aufzusuchen haben.

Es sei $p$ eine beliebige Primzahl (auch $p = 2$ nicht ausgeschlossen) und
\weq{1}{D' = p^2 D.}
Es sei $-D = m$, $-D' = m'$, also $m' = p^2 m$, und
\weq{2}{H_m(u) = 0,}
\weq{3}{H_{m'}(v) = 0}
seien die zu den Diskriminanten $-m, -m'$ gehörigen Klassengleichungen.
\[
v = j(\omega')
\]
sei eine beliebige Wurzel der zweiten und
\weq{4}{A'\omega'^2 + B'\omega' + C' = 0, \quad B'^2 - 4A'C' = D'}
die primitive quadratische Gleichung, der $\omega'$ genügt. $A'$ möge, was erlaubt ist, durch $p$ unteilbar vorausgesetzt sein.

Wir betrachten nun die Invariantengleichung
\weq{5}{F_p(u, v) = 0,}
deren Wurzeln sind:
\weq{6}{u = j(p\omega'), \quad j\left(\frac{\omega' + c}{p}\right), \quad c = 0, 1, 2, \ldots, p - 1.}
Das Argument dieser Funktionen:
\[
\omega = p\omega' \quad \text{oder} \quad = \frac{\omega' + c}{p}
\]
genügt einer aus (4) abzuleitenden quadratischen Gleichung, nämlich:
\weq{7}{\begin{gathered}
\omega = p\omega', \quad A'\omega^2 + B'p\omega + C'p^2 = 0, \\
\omega = \frac{\omega' + c}{p}, \\
A'p^2\omega^2 + (B' - 2A'c)p\omega + A'c^2 - B'c + C' = 0.
\end{gathered}}
Die Diskriminanten dieser Gleichungen sind $p^2D' = p^4D$, und die erste von ihnen ist immer primitiv, die Diskriminante
%CONT>>
